\documentclass[12pt,a4paper]{article}
\usepackage[margin=25mm, headheight=15pt, headsep=10pt]{geometry}
\usepackage{fontspec}
\usepackage[table]{xcolor} % Note: table option is required for rowcolors
\usepackage{titlesec}
\usepackage{tocloft}
\usepackage{fancyhdr}
\usepackage{booktabs}
\usepackage{tcolorbox}
\tcbuselibrary{skins, breakable}
\usepackage[colorlinks=true, linkcolor=emerald, urlcolor=emerald, bookmarks=true]{hyperref}

\definecolor{emerald}{RGB}{0,102,51}
\definecolor{darkred}{RGB}{139,0,0}
\definecolor{lightgray}{RGB}{245,245,245}

\usepackage[nil]{babel}
\babelprovide[import, main]{bengali}
\babelprovide[import]{arabic}
\babelfont{rm}[Renderer=HarfBuzz,Scale=1.0]{Noto Serif Bengali}
\babelfont{sf}[Renderer=HarfBuzz]{Noto Sans Bengali}
\babelfont{tt}[Renderer=HarfBuzz]{Noto Sans Bengali}
\babelfont[arabic]{rm}[Renderer=HarfBuzz,Scale=1.1]{Noto Naskh Arabic}

% Section styling
\titleformat{\section}{\Large\bfseries\color{emerald}}{\thesection.}{0.5em}{}
\titleformat{\subsection}{\large\bfseries\color{emerald!85!black}}{\thesubsection.}{0.5em}{}
\titleformat{\subsubsection}{\normalsize\bfseries\color{emerald!70!black}}{\thesubsubsection.}{0.5em}{}
\titlespacing*{\section}{0pt}{18pt}{8pt}

% Fancy Header/Footer setup
\pagestyle{fancy}
\fancyhf{} % clear all header and footer fields
\fancyhead[L]{\color{emerald}\small\textbf{\nouppercase{\leftmark}}}
\fancyfoot[L]{\scriptsize\color{black!50}{\lat{AI}}-সংকলিত, আলেম-যাচাইকৃত নয়}
\fancyfoot[C]{\color{emerald!70!black}\thepage}
\renewcommand{\headrulewidth}{0.4pt}
\renewcommand{\headrule}{\hbox to\headwidth{\color{emerald}\leaders\hrule height \headrulewidth\hfill}}
\renewcommand{\footrulewidth}{0pt}

% Custom boxes
% 1. Ayat/Hadith Box
\newtcolorbox{ayatbox}[1][]{
    enhanced, breakable, 
    colback=emerald!5, 
    colframe=emerald, 
    boxrule=0.5pt, 
    arc=3pt, 
    left=10pt, right=10pt, top=10pt, bottom=10pt,
    #1
}

% 2. Quote/Translation Box
\newtcolorbox{quotebox}[1][]{
    enhanced, breakable,
    frame hidden,
    colback=lightgray,
    borderline west={3pt}{0pt}{emerald},
    left=15pt, right=10pt, top=10pt, bottom=10pt,
    sharp corners,
    #1
}

% 3. AI-generated notice box
\newtcolorbox{warnbox}[1][]{
    enhanced, breakable,
    colback=red!4,
    colframe=darkred,
    boxrule=0.8pt,
    arc=3pt,
    left=10pt, right=10pt, top=10pt, bottom=10pt,
    #1
}

% Commands
\newcommand{\arab}[1]{\foreignlanguage{arabic}{#1}}
\newcommand{\kib}[1]{\textcolor{emerald}{\textbf{#1}}}

% Latin (English) fallback: Noto Serif Bengali has NO Latin glyphs
\newfontfamily\latfont[Renderer=HarfBuzz]{Noto Serif}
\newcommand{\lat}[1]{{\latfont #1}}

\setlength{\parskip}{5pt}
\setlength{\parindent}{0pt}

% Fix for missing bullet in Noto Serif Bengali
\renewcommand{\labelitemi}{$\bullet$}



\begin{document}
\begin{center}{\Large\bfseries\color{emerald} আল্লাহর রহমত ও তাওবা — হাদিস ও আসার}\end{center}
\vspace{1em}
\section*{হাদিস ও আসার — আল্লাহর রহমত, ক্ষমা ও তাওবা}

\begin{warnbox}
 \textbf{গুরুত্বপূর্ণ নোটিশ (\lat{Important} \lat{Notice}):} এই কন্টেন্টটি \lat{AI} (কৃত্রিম বুদ্ধিমত্তা) দিয়ে প্রস্তুত ও সংকলিত। \lat{AI} কঠোর সূত্র-মান নীতি মেনে শুধু নির্ভরযোগ্য উৎস থেকে তথ্য সংগ্রহ করেছে — কেবল সহীহ/হাসান হাদিস ও সহীহ সনদের আসার রাখা হয়েছে, দুর্বল সনদ বাদ দেওয়া হয়েছে। তবে এটি কোনো আলেম বা মুহাদ্দিস কর্তৃক যাচাইকৃত নয়।
\end{warnbox}

\begin{quote}
\textbf{পড়ার নিয়ম:} প্রতিটি হাদিসের সাথে — \textbf{সূত্র কার্ড} (সনদের মান) $\rightarrow$ \textbf{পটভূমি} $\rightarrow$ \textbf{পূর্ণ বর্ণনা (বাংলা, \lat{pure})} $\rightarrow$ \textbf{শব্দের ব্যাখ্যা} $\rightarrow$ \textbf{গুরুত্বপূর্ণ শিক্ষা}। \\
\textbf{মূল নীতিমালা:} শুধুমাত্র সহীহ/হাসান হাদিস। দুর্বল (যয়ীফ) বর্ণনা \textbf{দলিল হিসেবে নয়}; প্রচলিত-কিন্তু-দুর্বল বিষয়গুলো শেষের \lat{“}যাচাই টেবিলে\lat{”} স্পষ্টভাবে চিহ্নিত করা হয়েছে।
\end{quote}

\medskip\hrule\medskip

\subsection*{পার্ট ১ — আল্লাহর রহমতের ব্যাপ্তি (\lat{The} \lat{Vastness} \lat{of} \lat{Mercy})}

\subsubsection*{হাদিস ১ — \lat{“}আমার রহমত আমার গজবের উপর প্রাধান্য পায়\lat{”} (সহীহ মুসলিম ২৭৫১)}

\subsubsection*{১. সূত্র কার্ড}

\begin{table}[h]\centering\small
\begin{tabular}{ll}
বিষয় & বিবরণ \\
\hline
\textbf{বর্ণনাকারী} & আবু হুরাইরা (রাঃ) \\
\textbf{গ্রন্থ} & সহীহ মুসলিম ২৭৫১; সহীহ বুখারি ৩১৯৪, ৭৪৫৩ \\
\textbf{অধ্যায়} & কিতাবুত তাওবাহ — রহমতের ব্যাপকতা ও তা গজবের উপর বিজয়ী \\
\textbf{সনদের মান} & \textbf{সহীহ} (বুখারি-মুসলিম) \\
\hline
\end{tabular}
\end{table}

\subsubsection*{২. পটভূমি}

আল্লাহ সৃষ্টি শুরু করার সময় নিজের কিতাবে লিখে রেখেছেন — রহমত গজবের উপর প্রাধান্য পাবে। এটি রহমতের \textbf{সাংবিধানিক (\lat{constitutional})} ঘোষণা — সৃষ্টির আগে থেকেই রহমতের বিজয় স্থির।

\subsubsection*{৩. পূর্ণ বর্ণনা (বাংলা, \lat{pure})}

\begin{quote}
\textbf{\lat{“}আল্লাহ যখন সৃষ্টি সৃষ্টি করলেন, তখন তিনি তাঁর কিতাবে লিখলেন — আর তা তাঁর কাছে আরশের উপর: নিশ্চয়ই আমার রহমত আমার গজবের উপর বিজয়ী।\lat{”}}
\end{quote}

\subsubsection*{৪. শব্দের ব্যাখ্যা}

\begin{table}[h]\centering\small
\begin{tabular}{ll}
শব্দ & অর্থ \\
\hline
\textbf{তাগলিবু (\arab{تغلب})} & বিজয়ী/প্রাধান্য পায় \\
\textbf{কিতাব} & লেখা — আল্লাহর পূর্ব-নির্ধারিত ব্যবস্থা \\
\hline
\end{tabular}
\end{table}

\subsubsection*{৫. গুরুত্বপূর্ণ শিক্ষা}

\begin{enumerate}
\item \textbf{রহমত গজবের আগে:} আল্লাহর দয়াই তাঁর সৃষ্টির মূল ব্যবস্থা; শাস্তি ব্যতিক্রম, রহমত মূলনীতি।
\item \textbf{পাপীর আশা:} এই হাদিসের আলোকে — গুনাহ যতই বড় হোক, রহমতের ব্যবস্থা তার চেয়ে আগে লেখা।
\item \textbf{ভারসাম্য:} রহমতের বিজয় মানে পাপে নিরাপত্তা নয় — তাওবার পথ ধরলেই রহমত প্রযোজ্য।
\end{enumerate}
\medskip\hrule\medskip

\subsubsection*{হাদিস ২ — রহমতের একশো ভাগ (সহীহ মুসলিম ২৭৫২)}

\subsubsection*{১. সূত্র কার্ড}

\begin{table}[h]\centering\small
\begin{tabular}{ll}
বিষয় & বিবরণ \\
\hline
\textbf{বর্ণনাকারী} & আবু হুরাইরা (রাঃ) \\
\textbf{গ্রন্থ} & সহীহ মুসলিম ২৭৫২; সহীহ বুখারি ৬০০০ \\
\textbf{অধ্যায়} & কিতাবুত তাওবাহ — রহমতের ব্যাপকতা \\
\textbf{সনদের মান} & \textbf{সহীহ} (বুখারি-মুসলিম) \\
\hline
\end{tabular}
\end{table}

\subsubsection*{২. পূর্ণ বর্ণনা (বাংলা, \lat{pure})}

\begin{quote}
\textbf{\lat{“}আল্লাহ রহমতকে একশো ভাগে ভাগ করেছেন। তিনি তাঁর কাছে নিরানব্বই ভাগ রেখেছেন এবং পৃথিবীতে এক ভাগ নাজিল করেছেন। এই এক ভাগের কারণেই সমস্ত সৃষ্টি পরস্পরের প্রতি দয়া করে — এমনকি পশুও তার বাচ্চার উপর পা তোলার ভয়ে পা তুলে নেয়।\lat{”}}
\end{quote}

\subsubsection*{৩. গুরুত্বপূর্ণ শিক্ষা}

\begin{enumerate}
\item \textbf{পৃথিবীর সব দয়া = এক ভাগ:} মানুষের মধ্যে যে সামান্য মায়া-দয়া দেখা যায়, তা রহমতের মাত্র এক ভাগ — বাকি নিরানব্বই ভাগ আল্লাহর কাছে, কিয়ামতে।
\item \textbf{ক্ষমার উৎস:} আল্লাহর ক্ষমা-দয়ার ভাণ্ডার পৃথিবীর সব দয়ার চেয়ে বিশাল — তাই পাপীর আশা কখনো ফুরায় না।
\item \textbf{মায়ের স্নেহের প্রমাণ:} এই এক ভাগ দয়াতেই মা-বাবা সন্তানকে এত ভালোবাসেন; আল্লাহর দয়া তার চেয়ে অসীম (পরের হাদিস)।
\end{enumerate}
\medskip\hrule\medskip

\subsubsection*{হাদিস ৩ — আল্লাহ মায়ের চেয়েও বেশি দয়ালু (সহীহ বুখারি ৫৯৯৯, সহীহ মুসলিম ২৭৫৪)}

\subsubsection*{১. সূত্র কার্ড}

\begin{table}[h]\centering\small
\begin{tabular}{ll}
বিষয় & বিবরণ \\
\hline
\textbf{বর্ণনাকারী} & উমার ইবনে খাত্তাব (রাঃ) \\
\textbf{গ্রন্থ} & সহীহ বুখারি ৫৯৯৯; সহীহ মুসলিম ২৭৫৪ \\
\textbf{অধ্যায়} & কিতাবুল আদাব — রহমত \\
\textbf{সনদের মান} & \textbf{সহীহ} (বুখারি-মুসলিম) \\
\hline
\end{tabular}
\end{table}

\subsubsection*{২. পটভূমি}

যুদ্ধবন্দীদের মধ্যে একজন মহিলা তার শিশুকে খুঁজছিলেন — সন্তান পেয়ে তাকে জড়িয়ে ধরলেন। নবীজি (সা.) সাহাবিদের জিজ্ঞেস করলেন — এই মহিলা কি স্বেচ্ছায় নিজের সন্তানকে আগুনে ফেলবে? সবাই বললেন — না। তখন নবীজি (সা.) বললেন — আল্লাহ তাঁর বান্দাদের প্রতি এই মহিলার চেয়েও বেশি দয়ালু।

\subsubsection*{৩. পূর্ণ বর্ণনা (বাংলা, \lat{pure})}

\begin{quote}
\textbf{\lat{“}আল্লাহ তাঁর বান্দাদের প্রতি এই (মহিলার) চেয়েও বেশি দয়ালু।\lat{”}}
\end{quote}

\subsubsection*{৪. গুরুত্বপূর্ণ শিক্ষা}

\begin{enumerate}
\item \textbf{উপমার গভীরতা:} মায়ের সন্তানের প্রতি ভালোবাসা মানুষের মধ্যে সবচেয়ে শক্তিশালী দয়া — আল্লাহর দয়া তার চেয়েও বেশি।
\item \textbf{পাপীর জন্য বার্তা:} মানুষ নিজের সন্তানকে ক্ষমা করে; আল্লাহ পাপী বান্দাকে তার চেয়ে বেশি ভালোবাসেন — যখন সে ফিরে আসে।
\end{enumerate}
\medskip\hrule\medskip

\subsubsection*{হাদিস ৪ — হাদিস কুদসি: \lat{“}হে আমার বান্দারা, তোমরা রাতে-দিনে পাপ করো…\lat{”} (সহীহ মুসলিম ২৫৭৭)}

\subsubsection*{১. সূত্র কার্ড}

\begin{table}[h]\centering\small
\begin{tabular}{ll}
বিষয় & বিবরণ \\
\hline
\textbf{বর্ণনাকারী} & আবু যার গিফারী (রাঃ) \\
\textbf{গ্রন্থ} & সহীহ মুসলিম ২৫৭৭ \\
\textbf{অধ্যায়} & কিতাবুল বিরর — জুলুম নিষেধ \\
\textbf{সনদের মান} & \textbf{সহীহ} \\
\hline
\end{tabular}
\end{table}

\subsubsection*{২. পূর্ণ বর্ণনা (বাংলা, \lat{pure})}

\begin{quote}
আল্লাহ বলেন: \textbf{\lat{“}হে আমার বান্দারা! আমি নিজের জন্য জুলুম হারাম করেছি এবং তোমাদের মধ্যে তা হারাম বানিয়েছি — সুতরাং তোমরা একে অপরের উপর জুলুম করো না… হে আমার বান্দারা! তোমরা রাতে ও দিনে পাপ করতে থাকো, আর আমি সব পাপ ক্ষমা করি — সুতরাং তোমরা আমার কাছে ক্ষমা চাও, আমি তোমাদের ক্ষমা করব।\lat{”}}
\end{quote}

\subsubsection*{৩. শব্দের ব্যাখ্যা}

\begin{table}[h]\centering\small
\begin{tabular}{ll}
শব্দ & অর্থ \\
\hline
\textbf{তুনযিবূনা (\arab{تذنبون})} & পাপ করতে থাকো \\
\textbf{আগফিরুয যুনূবা জামীআন} & সব পাপ ক্ষমা করি \\
\textbf{ফাসতাগফিরূনী} & সুতরাং আমার কাছে ক্ষমা চাও \\
\hline
\end{tabular}
\end{table}

\subsubsection*{৪. গুরুত্বপূর্ণ শিক্ষা}

\begin{enumerate}
\item \textbf{আল্লাহ নিজেই স্বীকার করেছেন:} বান্দারা পাপ করে — আল্লাহ বলেন \lat{“}তোমরা পাপ করো, আমি ক্ষমা করি\lat{”} — অর্থ পাপীর জন্য দরজা খোলা রাখা আল্লাহর নিজের ঘোষণা।
\item \textbf{পাপের পর করণীয়:} ক্ষমা চাওয়া — শুধু এটাই পাপের সঠিক উত্তর; হতাশা নয়, আত্মঘাতিতা নয়।
\item \textbf{জুলুমের ব্যতিক্রম:} একই হাদিসে মানুষের হকের জুলুম নিষিদ্ধ — তাওবার সাথে মানুষের হক আদায়ও জরুরি।
\end{enumerate}
\medskip\hrule\medskip

\subsubsection*{হাদিস ৫ — \lat{“}কেউ নিজের আমলে জান্নাতে যাবে না\lat{”} (সহীহ বুখারি ৫৬৭৩, সহীহ মুসলিম ২৮১৬)}

\subsubsection*{১. সূত্র কার্ড}

\begin{table}[h]\centering\small
\begin{tabular}{ll}
বিষয় & বিবরণ \\
\hline
\textbf{বর্ণনাকারী} & আবু হুরাইরা (রাঃ) \\
\textbf{গ্রন্থ} & সহীহ বুখারি ৫৬৭৩; সহীহ মুসলিম ২৮১৬ \\
\textbf{অধ্যায়} & কিতাবুর রিকাক — আমলের উপর নির্ভর না করে রহমতের আশা \\
\textbf{সনদের মান} & \textbf{সহীহ} (বুখারি-মুসলিম) \\
\hline
\end{tabular}
\end{table}

\subsubsection*{২. পূর্ণ বর্ণনা (বাংলা, \lat{pure})}

\begin{quote}
নবীজি (সা.) বললেন: \textbf{\lat{“}তোমাদের কেউ তার আমল দ্বারা কখনো (জান্নাতে) প্রবেশ করবে না।\lat{”}} সাহাবিরা জিজ্ঞেস করলেন: হে আল্লাহর রাসূল! আপনিও না? তিনি বললেন: \textbf{\lat{“}আমিও না — তবে আল্লাহ যদি আমাকে তাঁর রহমত ও অনুগ্রহে আবৃত করেন।\lat{”}}
\end{quote}

\subsubsection*{৩. গুরুত্বপূর্ণ শিক্ষা}

\begin{enumerate}
\item \textbf{নবীজিও রহমতের মুখাপেক্ষী:} আমল যত ভালোই হোক, চূড়ান্ত নির্ভরতা আল্লাহর রহমতের উপর — এটাই মুমিনের বিনয়।
\item \textbf{পাপীর আশা:} রহমত ছাড়া কারো গতি নেই; তাই পাপী ব্যক্তি রহমতের আশা করতেই পারে — রহমতই সবার পথ।
\item \textbf{আমল ত্যাগের অজুহাত নয়:} আমলই একমাত্র পথ নয়, কিন্তু আমল ছাড়াও নয় — আমল রহমতের দিকে অগ্রসর হওয়ার উপায়।
\end{enumerate}
\medskip\hrule\medskip

\subsection*{পার্ট ২ — তাওবার দরজা খোলা (\lat{The} \lat{Door} \lat{is} \lat{Open})}

\subsubsection*{হাদিস ৬ — আল্লাহ রাতে-দিনে হাত বাড়ান (সহীহ মুসলিম ২৭৫৯)}

\subsubsection*{১. সূত্র কার্ড}

\begin{table}[h]\centering\small
\begin{tabular}{ll}
বিষয় & বিবরণ \\
\hline
\textbf{বর্ণনাকারী} & আবু মুসা আশ'আরি (রাঃ) \\
\textbf{গ্রন্থ} & সহীহ মুসলিম ২৭৫৯ \\
\textbf{অধ্যায়} & কিতাবুত তাওবাহ — পাপ ও তাওবা বারবার হলেও তাওবা কবুল \\
\textbf{সনদের মান} & \textbf{সহীহ} \\
\hline
\end{tabular}
\end{table}

\subsubsection*{২. পূর্ণ বর্ণনা (বাংলা, \lat{pure})}

\begin{quote}
\textbf{\lat{“}নিশ্চয়ই আল্লাহ রাতে তাঁর হাত বাড়ান, যাতে দিনের পাপী তাওবা করে; আর দিনে তাঁর হাত বাড়ান, যাতে রাতের পাপী তাওবা করে — এমনকি পশ্চিম দিক থেকে সূর্য উদিত হওয়ার পূর্ব পর্যন্ত।\lat{”}}
\end{quote}

\subsubsection*{৩. শব্দের ব্যাখ্যা}

\begin{table}[h]\centering\small
\begin{tabular}{ll}
শব্দ & অর্থ \\
\hline
\textbf{ইয়াবসুতু ইয়াদাহু (\arab{يبسط} \arab{يده})} & তাঁর হাত বাড়ান — তাওবা গ্রহণের প্রস্তুতি \\
\textbf{মুসিউন নাহার (\arab{مسيء} \arab{النهار})} & দিনের পাপী \\
\textbf{হাত্তা ত্বালউআশ শামসু} & যতক্ষণ না সূর্য পশ্চিম থেকে উদিত হয় \\
\hline
\end{tabular}
\end{table}

\subsubsection*{৪. গুরুত্বপূর্ণ শিক্ষা}

\begin{enumerate}
\item \textbf{সময় সীমা:} তাওবার দরজা খোলা — মৃত্যু বা পশ্চিমে সূর্যোদয় পর্যন্ত; এখনই সময়, কিন্তু সময় আছে।
\item \textbf{বারবার পাপ-তাওবা:} অধ্যায়ের নামই — \lat{“}পাপ ও তাওবা বারবার হলেও কবুল\lat{”} — অর্থ: আবার পাপ করলে আবার তাওবা; দরজা বন্ধ হয় না।
\item \textbf{আল্লাহর পক্ষ থেকে উদ্যোগ:} হাত বাড়ানো আল্লাহর পক্ষ থেকে — বান্দার দিকে আগে এগিয়ে আসা তাঁরই রীতি (মুসলিম ২৬৭৫ দেখুন)।
\end{enumerate}
\medskip\hrule\medskip

\subsubsection*{হাদিস ৭ — প্রাণ কণ্ঠাগত হওয়ার আগ পর্যন্ত তাওবা কবুল (তিরমিযী ৩৫৩৭)}

\subsubsection*{১. সূত্র কার্ড}

\begin{table}[h]\centering\small
\begin{tabular}{ll}
বিষয় & বিবরণ \\
\hline
\textbf{বর্ণনাকারী} & আবদুল্লাহ ইবনে উমার (রাঃ) \\
\textbf{গ্রন্থ} & জামি' আত-তিরমিযী ৩৫৩৭ \\
\textbf{অধ্যায়} & কিতাবুদ দাওয়াত — আল্লাহ তাওবা কবুল করেন যতক্ষণ প্রাণ কণ্ঠাগত না হয় \\
\textbf{সনদের মান} & \textbf{হাসান} (তিরমিযী: হাসান-গারীব; দারুসসালাম: হাসান) \\
\hline
\end{tabular}
\end{table}

\subsubsection*{২. পূর্ণ বর্ণনা (বাংলা, \lat{pure})}

\begin{quote}
\textbf{\lat{“}নিশ্চয়ই আল্লাহ বান্দার তাওবা কবুল করেন — যতক্ষণ না তার প্রাণ কণ্ঠাগত হয় (মৃত্যু নিকটবর্তী হয়)।\lat{”}}
\end{quote}

\subsubsection*{৩. শব্দের ব্যাখ্যা}

\begin{table}[h]\centering\small
\begin{tabular}{ll}
শব্দ & অর্থ \\
\hline
\textbf{ইয়ুগারগির (\arab{يغرغر})} & প্রাণ কণ্ঠাগত হওয়া — মৃত্যুর শেষ মুহূর্ত \\
\hline
\end{tabular}
\end{table}

\subsubsection*{৪. গুরুত্বপূর্ণ শিক্ষা}

\begin{enumerate}
\item \textbf{দরজা শেষ মুহূর্ত পর্যন্ত খোলা:} মৃত্যুর আগমনের পূর্ব পর্যন্ত তাওবা কবুল — পাপী যত দেরিই করুক, দরজা তখনো খোলা (তবে দেরি করা বোকামি — মৃত্যু কখন আসে কেউ জানে না)।
\item \textbf{হতাশার জবাব:} \lat{“}এত দেরি হয়ে গেছে\lat{”} — এই চিন্তার জবাব এই হাদিস: যতক্ষণ শ্বাস আছে, দরজা আছে।
\end{enumerate}
\medskip\hrule\medskip

\subsubsection*{হাদিস ৮ — তাওবায় আল্লাহর খুশি: মরুভূমিতে হারানো উট (সহীহ বুখারি ৬৩০৮, সহীহ মুসলিম ২৭৪৪)}

\subsubsection*{১. সূত্র কার্ড}

\begin{table}[h]\centering\small
\begin{tabular}{ll}
বিষয় & বিবরণ \\
\hline
\textbf{বর্ণনাকারী} & আবদুল্লাহ ইবনে মাসউদ (রাঃ) \\
\textbf{গ্রন্থ} & সহীহ বুখারি ৬৩০৮; সহীহ মুসলিম ২৭৪৪ \\
\textbf{অধ্যায়} & বুখারি: কিতাবুদ দাওয়াত — বাবুত তাওবাহ \\
\textbf{সনদের মান} & \textbf{সহীহ} (বুখারি-মুসলিম) \\
\hline
\end{tabular}
\end{table}

\subsubsection*{২. পটভূমি}

ইবনে মাসউদ (রাঃ) দুইটি বর্ণনা দিয়েছেন — একটি নবীজির (সা.) থেকে, আরেকটি নিজের থেকে। নিজের বর্ণনায় তিনি বলেন — মুমিন তার পাপকে এমন দেখে, যেন পাহাড়ের নিচে বসে আছে, যা যে-কোনো মুহূর্তে তার উপর ভেঙে পড়বে; আর পাপী তার পাপকে মাছির মতো দেখে, নাকের উপর উড়ে গেলেই দূর করে দেয়।

\subsubsection*{৩. পূর্ণ বর্ণনা (বাংলা, \lat{pure})}

\begin{quote}
নবীজি (সা.) বলেছেন: \textbf{\lat{“}আল্লাহ বান্দার তাওবায় এত বেশি খুশি হন, যত বেশি খুশি হয় তোমাদের কেউ — যখন সে এমন এক মরুভূমিতে তার উট হারায় (যেখানে খাবার-পানি ছিল), সেটা খুঁজে না পেয়ে মৃত্যুর অপেক্ষায় ঘুমিয়ে পড়ে, তারপর ঘুম থেকে উঠে দেখে তার উটটি তার সামনে দাঁড়িয়ে!\lat{”}}
\end{quote}

\subsubsection*{৪. শব্দের ব্যাখ্যা}

\begin{table}[h]\centering\small
\begin{tabular}{ll}
শব্দ & অর্থ \\
\hline
\textbf{আশাদ্দু ফারাহান (\arab{أشد} \arab{فرحا})} & সবচেয়ে বেশি খুশি \\
\textbf{আদ-দাল্লাহ (\arab{الضالة})} & হারানো উট/জিনিস \\
\textbf{আল-ফালাহ (\arab{الفلاة})} & মরুভূমি — প্রাণহীন এলাকা \\
\hline
\end{tabular}
\end{table}

\subsubsection*{৫. গুরুত্বপূর্ণ শিক্ষা}

\begin{enumerate}
\item \textbf{রহমতের দৃষ্টান্ত:} মরুভূমিতে মৃত্যুর মুখ থেকে উট ফিরে পাওয়ার আনন্দ — আল্লাহর কাছে তাওবাকারী বান্দার প্রত্যাবর্তন এর চেয়ে আনন্দের।
\item \textbf{পাপের গুরুত্ব:} ইবনে মাসউদের (রাঃ) নিজের বর্ণনা — মুমিন পাপকে পাহাড়ের মতো ভয় পায় (পাপের ভয় থাকে), কিন্তু হতাশ হয় না (তাওবা করে)। পাপকে মাছির মতো হালকা ভাবাই আসল বিপদ।
\item \textbf{তাওবার আগে আল্লাহর খুশি:} আল্লাহ তাওবার \lat{“}ফল\lat{”} নয় — তাওবার \lat{“}আগেই\lat{”} খুশি হন; অর্থাৎ আল্লাহ চান বান্দা ফিরুক।
\end{enumerate}
\medskip\hrule\medskip

\subsubsection*{হাদিস ৯ — \lat{“}আমি আমার বান্দার ধারণা অনুযায়ী…\lat{”} (সহীহ মুসলিম ২৬৭৫)}

\subsubsection*{১. সূত্র কার্ড}

\begin{table}[h]\centering\small
\begin{tabular}{ll}
বিষয় & বিবরণ \\
\hline
\textbf{বর্ণনাকারী} & আবু হুরাইরা (রাঃ) \\
\textbf{গ্রন্থ} & সহীহ মুসলিম ২৬৭৫; সহীহ বুখারি ৭৪০৫ \\
\textbf{অধ্যায়} & কিতাবুত তাওবাহ — তাওবার উৎসাহ \\
\textbf{সনদের মান} & \textbf{সহীহ} (বুখারি-মুসলিম) \\
\hline
\end{tabular}
\end{table}

\subsubsection*{২. পূর্ণ বর্ণনা (বাংলা, \lat{pure})}

\begin{quote}
আল্লাহ বলেন: \textbf{\lat{“}আমি আমার বান্দার আমার সম্পর্কে ধারণা অনুযায়ী থাকি — আর সে যখন আমাকে স্মরণ করে, আমি তার সাথে থাকি।\lat{”}} অতঃপর নবীজি (সা.) বললেন: \textbf{\lat{“}আল্লাহর কসম! আল্লাহ বান্দার তাওবায় এত বেশি খুশি হন, যত বেশি খুশি হন তোমাদের কেউ মরুভূমিতে তার হারানো উট ফিরে পেলে। যে আমার দিকে এক বিঘত এগোয়, আমি তার দিকে এক হাত এগিয়ে যাই; যে এক হাত এগোয়, আমি এক বাহু এগিয়ে যাই; আর যে আমার দিকে হেঁটে আসে, আমি তার দিকে দৌড়ে যাই।\lat{”}}
\end{quote}

\subsubsection*{৩. শব্দের ব্যাখ্যা}

\begin{table}[h]\centering\small
\begin{tabular}{ll}
শব্দ & অর্থ \\
\hline
\textbf{আনা ইনদা যান্নি আবদি বী} & আমি আমার বান্দার আমার সম্পর্কে ধারণা অনুযায়ী \\
\textbf{শিবরান (\arab{شبرا})} & এক বিঘত (হাতের প্রস্থ) \\
\textbf{উহারওইলু (\arab{أهرول})} & দ্রুত/দৌড়ে এগিয়ে যাওয়া \\
\hline
\end{tabular}
\end{table}

\subsubsection*{৪. গুরুত্বপূর্ণ শিক্ষা}

\begin{enumerate}
\item \textbf{ধারণা গুরুত্বপূর্ণ:} বান্দা আল্লাহকে ক্ষমাশীল ভাবলে — সে ক্ষমা পাবে; \lat{“}আল্লাহ আমাকে ক্ষমা করবেন না\lat{”} ভাবলে সেই ধারণাই তাকে আটকে রাখে (আল্লাহর ক্ষমা সম্পর্কে খারাপ ধারণা পোষণ করা নিষিদ্ধ — যুমার ৩৯:৫৩)।
\item \textbf{নৈকট্যের আনুপাতিকতা:} সামান্য এগোলেও আল্লাহ বেশি এগিয়ে আসেন — তাওবার প্রথম পদক্ষেপই বিশাল।
\item \textbf{স্মরণের ফল:} \lat{“}স্মরণ করলে আমি তার সাথে থাকি\lat{”} — তাওবাকারী বান্দার সাথে আল্লাহর বিশেষ উপস্থিতি।
\end{enumerate}
\medskip\hrule\medskip

\subsubsection*{হাদিস ১০ — পাপ না করলে আল্লাহ অন্য জাতি আনতেন (সহীহ মুসলিম ২৭৪৯)}

\subsubsection*{১. সূত্র কার্ড}

\begin{table}[h]\centering\small
\begin{tabular}{ll}
বিষয় & বিবরণ \\
\hline
\textbf{বর্ণনাকারী} & আবু হুরাইরা (রাঃ) \\
\textbf{গ্রন্থ} & সহীহ মুসলিম ২৭৪৯ \\
\textbf{অধ্যায়} & কিতাবুত তাওবাহ — ইস্তিগফার ও তাওবায় পাপ মুছে যায় \\
\textbf{সনদের মান} & \textbf{সহীহ} \\
\hline
\end{tabular}
\end{table}

\subsubsection*{২. পূর্ণ বর্ণনা (বাংলা, \lat{pure})}

\begin{quote}
\textbf{\lat{“}যাঁর হাতে আমার প্রাণ, তাঁর কসম! তোমরা যদি পাপ না করতে, তবে আল্লাহ তোমাদের সরিয়ে দিয়ে এমন এক জাতি আনতেন, যারা পাপ করত এবং আল্লাহর কাছে ক্ষমা চাইত — আর আল্লাহ তাদের ক্ষমা করে দিতেন।\lat{”}}
\end{quote}

\subsubsection*{৩. গুরুত্বপূর্ণ শিক্ষা}

\begin{enumerate}
\item \textbf{পাপ-ক্ষমার চক্রই আল্লাহর নকশা:} এই দুনিয়ায় \lat{“}পাপ $\rightarrow$ ইস্তিগফার $\rightarrow$ ক্ষমা\lat{”} — এই সম্পর্কই আল্লাহ চান; পাপী না থাকলে ক্ষমার দরজা বন্ধ থাকত।
\item \textbf{পাপ অজুহাত নয়, তাওবা অজুহাত:} হাদিস পাপের প্রশংসা করে না — প্রশংসা করে পাপের পরের ইস্তিগফারকে; পাপে ডুবে থাকা নয়, পাপ থেকে ফেরা।
\item \textbf{অসম্ভব লক্ষ্যের চাপ নেই:} পাপশূন্য জীবন কারো পক্ষে সম্ভব না — তাই পাপ হলে ভেঙে না পড়ে তাওবার দিকে ছুটতে হবে।
\end{enumerate}
\medskip\hrule\medskip

\subsubsection*{হাদিস ১১ — \lat{“}তোমার গুনাহ আকাশের মেঘ পর্যন্ত পৌঁছালেও…\lat{”} (তিরমিযী ৩৫৪০)}

\subsubsection*{১. সূত্র কার্ড}

\begin{table}[h]\centering\small
\begin{tabular}{ll}
বিষয় & বিবরণ \\
\hline
\textbf{বর্ণনাকারী} & আনাস ইবনে মালিক (রাঃ) \\
\textbf{গ্রন্থ} & জামি' আত-তিরমিযী ৩৫৪০; মুসনাদ আহমাদ; সুনানে ইবনে মাজাহ ৪২৫৭ \\
\textbf{অধ্যায়} & কিতাবুদ দাওয়াত — \lat{“}হে আদম সন্তান, যতক্ষণ তুমি আমাকে ডাকবে…\lat{”} \\
\textbf{সনদের মান} & \textbf{হাসান} (দারুসসালাম: হাসান; আলবানি: হাসান) \\
\hline
\end{tabular}
\end{table}

\subsubsection*{২. পূর্ণ বর্ণনা (বাংলা, \lat{pure})}

\begin{quote}
আল্লাহ বলেন: \textbf{\lat{“}হে আদম সন্তান! যতক্ষণ তুমি আমাকে ডাকবে এবং আমার কাছে আশা রাখবে, ততক্ষণ আমি তোমাকে ক্ষমা করব — তোমার থেকে যা-ই হয়ে থাকুক, আমি পরোয়া করি না। হে আদম সন্তান! তোমার গুনাহ যদি আকাশের মেঘ পর্যন্ত পৌঁছে যায়, তারপর তুমি আমার কাছে ক্ষমা চাও — আমি তোমাকে ক্ষমা করব, পরোয়া করব না। হে আদম সন্তান! তুমি যদি পৃথিবী পরিমাণ পাপ নিয়ে আমার কাছে আসো, আমার সাথে কিছু শরিক না করে — আমি তোমার কাছে পৃথিবী পরিমাণ ক্ষমা নিয়ে আসব।\lat{”}}
\end{quote}

\subsubsection*{৩. শব্দের ব্যাখ্যা}

\begin{table}[h]\centering\small
\begin{tabular}{ll}
শব্দ & অর্থ \\
\hline
\textbf{আনানুস সামা (\arab{عنان} \arab{السماء})} & আকাশের মেঘ — বিশাল উচ্চতা \\
\textbf{কিরাবিল আরদ (\arab{قراب} \arab{الأرض})} & পৃথিবীর সমান পরিমাণ — পুরো পৃথিবী ভরতি পাপ \\
\textbf{লা উবালি (\arab{لا} \arab{أبالي})} & পরোয়া করি না — ক্ষমায় ক্লান্তি নেই \\
\hline
\end{tabular}
\end{table}

\subsubsection*{৪. গুরুত্বপূর্ণ শিক্ষা}

\begin{enumerate}
\item \textbf{আকারের তুলনা:} পাপের আকার (আকাশ-পৃথিবী) ক্ষমার আকারের (পৃথিবী পরিমাণ) চেয়ে বড় নয় — রহমত সবসময় বড়।
\item \textbf{একমাত্র শর্ত:} \lat{“}আমার সাথে শরিক না করে\lat{”} — শিরকমুক্ত অবস্থায় তাওবা; বাকি সব পাপ ক্ষমার আওতায়।
\item \textbf{ডাকাই দরজা:} \lat{“}ডাকবে ও আশা রাখবে\lat{”} — দোয়া ও আশা দুই-ই তাওবার অংশ।
\end{enumerate}
\medskip\hrule\medskip

\subsubsection*{হাদিস ১২ — পুড়ে-ছাই হওয়ার ভয় ও ক্ষমা (সহীহ মুসলিম ২৭৫৬)}

\subsubsection*{১. সূত্র কার্ড}

\begin{table}[h]\centering\small
\begin{tabular}{ll}
বিষয় & বিবরণ \\
\hline
\textbf{বর্ণনাকারী} & আবু হুরাইরা (রাঃ) \\
\textbf{গ্রন্থ} & সহীহ মুসলিম ২৭৫৬; সহীহ বুখারি ৩৪৭৮ \\
\textbf{অধ্যায়} & কিতাবুত তাওবাহ — রহমতের ব্যাপকতা \\
\textbf{সনদের মান} & \textbf{সহীহ} (বুখারি-মুসলিম) \\
\hline
\end{tabular}
\end{table}

\subsubsection*{২. পূর্ণ বর্ণনা (বাংলা, \lat{pure})}

\begin{quote}
এক ব্যক্তি কোনো ভালো কাজই করেনি। মৃত্যুর সময় সে পরিবারকে বলল: \textbf{\lat{“}আমি মারা গেলে আমাকে পুড়িয়ে ফেলো, এরপর আমার ছাইয়ের অর্ধেক স্থলে ছড়িয়ে দাও, অর্ধেক সাগরে। আল্লাহর কসম! যদি আল্লাহ আমাকে ধরতে পারেন, তিনি আমাকে এমন শাস্তি দেবেন, যা দুনিয়ার কাউকে দেননি।\lat{”}} তিনি মারা গেলে তাই করা হলো। আল্লাহ স্থলকে আদেশ দিলেন — ছাই জড়ো করল, সাগরকে আদেশ দিলেন — তাও জড়ো করল। আল্লাহ জিজ্ঞেস করলেন: কেন এমন করলে? সে বলল: \textbf{\lat{“}হে আমার রব! আপনার ভয়ে — আর আপনি তো সবচেয়ে ভালো জানেন।\lat{”}} আল্লাহ তাকে ক্ষমা করে দিলেন।
\end{quote}

\subsubsection*{৩. শব্দের ব্যাখ্যা}

\begin{table}[h]\centering\small
\begin{tabular}{ll}
শব্দ & অর্থ \\
\hline
\textbf{হাররিকূহু (\arab{حرقوه})} & পুড়িয়ে ফেলো \\
\textbf{ইযরূ (\arab{اذروا})} & ছড়িয়ে দাও \\
\textbf{মিন খাশইয়াতিকা (\arab{من} \arab{خشيتك})} & আপনার ভয়ের কারণে \\
\hline
\end{tabular}
\end{table}

\subsubsection*{৪. গুরুত্বপূর্ণ শিক্ষা}

\begin{enumerate}
\item \textbf{ভয়ই উদ্ধার করল:} এক নেক কাজহীন ব্যক্তির পরকালের ভয় — আল্লাহর শাস্তির অন্তরের ভয় — তাকে ক্ষমার দিকে নিয়ে গেল।
\item \textbf{আল্লাহর ক্ষমার ব্যাপকতা:} যে আল্লাহর ভয়ে ভীত — এমনকি তার ভয়ের প্রকাশ (পুড়ে যাওয়ার চিন্তা) ভুল হলেও — আল্লাহ তার অন্তরের ইখলাস দেখে ক্ষমা করেন।
\item \textbf{নিরাশ নয়, ভীত:} হাদিসের শিক্ষা — অন্তরে আল্লাহর প্রতি ভয় থাকলেই আশা আছে; হতাশা নয়, ভয় তাকে আল্লাহর কাছে পেশ করল।
\end{enumerate}
\medskip\hrule\medskip

\subsubsection*{হাদিস ১৩ — \lat{“}আমার বান্দা জানে তার রব আছে…\lat{”} — বারবার পাপ, বারবার ক্ষমা (সহীহ বুখারি ৭৫০৭)}

\subsubsection*{১. সূত্র কার্ড}

\begin{table}[h]\centering\small
\begin{tabular}{ll}
বিষয় & বিবরণ \\
\hline
\textbf{বর্ণনাকারী} & আবু হুরাইরা (রাঃ) \\
\textbf{গ্রন্থ} & সহীহ বুখারি ৭৫০৭; সহীহ মুসলিম ২৭৫৮ \\
\textbf{অধ্যায়} & কিতাবুত তাওহীদ — আল্লাহর বাণী সম্পর্কিত \\
\textbf{সনদের মান} & \textbf{সহীহ} (বুখারি-মুসলিম) \\
\hline
\end{tabular}
\end{table}

\subsubsection*{২. পূর্ণ বর্ণনা (বাংলা, \lat{pure})}

\begin{quote}
\textbf{\lat{“}এক বান্দা পাপ করল এবং বলল: হে আমার রব! আমি পাপ করেছি — আমাকে ক্ষমা করো। তার রব বললেন: আমার বান্দা জেনেছে — তার এমন রব আছে, যিনি পাপ ক্ষমা করেন এবং (চাইলে) পাপের শাস্তিও দেন। আমি আমার বান্দাকে ক্ষমা করে দিলাম। তারপর সে কিছুদিন পর আবার পাপ করল… এভাবে তৃতীয়বারও। আল্লাহ বললেন: আমার বান্দা জেনেছে তার এমন রব আছে, যিনি পাপ ক্ষমা করেন — আমি আমার বান্দাকে ক্ষমা করে দিলাম।\lat{”}}
\end{quote}

\subsubsection*{৩. গুরুত্বপূর্ণ শিক্ষা}

\begin{enumerate}
\item \textbf{বারবার পাপ $\rightarrow$ বারবার ক্ষমা:} তিনবার পাপ, তিনবার ক্ষমা — আল্লাহ পাপের পুনরাবৃত্তিতে তাওবা বন্ধ করেন না।
\item \textbf{ক্ষমার কারণ:} বান্দার জ্ঞান — \lat{“}আমার রব ক্ষমা করেন\lat{”} — এই জ্ঞানই তাকে ডাকায়; অর্থাৎ আল্লাহর ক্ষমার আশা রাখা নিজেই ক্ষমার পথ।
\item \textbf{সতর্কতা:} হাদিস পাপে উৎসাহ দেয় না — উৎসাহ দেয় ক্ষমা চাওয়ার অভ্যাসে; আলেমরা বলেন — দরজা খোলা থাকা মানে গেটে বসে থাকা নয়, ভেতরে ঢোকার চেষ্টা।
\end{enumerate}
\medskip\hrule\medskip

\subsubsection*{হাদিস ১৪ — \lat{“}সেরা গুনাহকারী তাওবাকারী\lat{”} (সুনানে ইবনে মাজাহ ৪২৫১)}

\subsubsection*{১. সূত্র কার্ড}

\begin{table}[h]\centering\small
\begin{tabular}{ll}
বিষয় & বিবরণ \\
\hline
\textbf{বর্ণনাকারী} & আনাস ইবনে মালিক (রাঃ) \\
\textbf{গ্রন্থ} & সুনানে ইবনে মাজাহ ৪২৫১; জামি' আত-তিরমিযী ২৪৯৯; মুসনাদ আহমাদ \\
\textbf{অধ্যায়} & কিতাবুজ যুহদ — তাওবার অধ্যায় \\
\textbf{সনদের মান} & \textbf{হাসান} (দারুসসালাম: হাসান; তিরমিযীর এক সনদে দুর্বলতার মত আছে — তবে অর্থ বহু সূত্রে সমর্থিত) \\
\hline
\end{tabular}
\end{table}

\subsubsection*{২. পূর্ণ বর্ণনা (বাংলা, \lat{pure})}

\begin{quote}
\textbf{\lat{“}প্রতিটি আদম সন্তান গুনাহ করে, আর গুনাহকারীদের মধ্যে সেরা তারা, যারা তাওবা করে।\lat{”}}
\end{quote}

\subsubsection*{৩. শব্দের ব্যাখ্যা}

\begin{table}[h]\centering\small
\begin{tabular}{ll}
শব্দ & অর্থ \\
\hline
\textbf{খাত্তা (\arab{خطّاء})} & গুনাহকারী — পাপপ্রবণ \\
\textbf{খাইরুল খাত্তাঈন} & গুনাহকারীদের মধ্যে সেরা \\
\textbf{আত-তাওওয়াবূন (\arab{التوابون})} & তাওবাকারীগণ \\
\hline
\end{tabular}
\end{table}

\subsubsection*{৪. গুরুত্বপূর্ণ শিক্ষা}

\begin{enumerate}
\item \textbf{সবার স্বভাব:} পাপ করা মানুষের স্বভাব — হাদিস এটাকে স্বাভাবিক বলেছে; তাই পাপে চমকে ভেঙে পড়া নয়।
\item \textbf{শ্রেষ্ঠত্ব তাওবায়:} মানুষের মধ্যে শ্রেষ্ঠ — যারা পাপের পর তাওবা করে; তাওবাই পাপীকে সেরা বানিয়ে দেয়।
\item \textbf{বিঃদ্রঃ} — তিরমিযীর (২৪৯৯) এক সনদে দুর্বলতার কথা কেউ কেউ বলেছেন; তবে ইবনে মাজাহর সনদ হাসান এবং অর্থ অন্যান্য হাদিসে (যেমন: মুসলিম ২৭৪৯, ২৭৫৯) সমর্থিত।
\end{enumerate}
\medskip\hrule\medskip

\subsubsection*{হাদিস ১৫ — \lat{“}অনুতাপই তাওবা\lat{”} (সুনানে ইবনে মাজাহ ৪২৫২)}

\subsubsection*{১. সূত্র কার্ড}

\begin{table}[h]\centering\small
\begin{tabular}{ll}
বিষয় & বিবরণ \\
\hline
\textbf{বর্ণনাকারী} & আবদুল্লাহ ইবনে মাসউদ (রাঃ) \\
\textbf{গ্রন্থ} & সুনানে ইবনে মাজাহ ৪২৫২; মুসনাদ আহমাদ \\
\textbf{অধ্যায়} & কিতাবুজ যুহদ — তাওবার অধ্যায় \\
\textbf{সনদের মান} & \textbf{হাসান} (দারুসসালাম: হাসান) \\
\hline
\end{tabular}
\end{table}

\subsubsection*{২. পূর্ণ বর্ণনা (বাংলা, \lat{pure})}

\begin{quote}
\textbf{\lat{“}অনুতাপই তাওবা।\lat{”}}
\end{quote}

\subsubsection*{৩. শব্দের ব্যাখ্যা}

\begin{table}[h]\centering\small
\begin{tabular}{ll}
শব্দ & অর্থ \\
\hline
\textbf{আন-নাদম (\arab{الندم})} & অনুতাপ/আফসোস — পাপ করার দুঃখ \\
\hline
\end{tabular}
\end{table}

\subsubsection*{৪. গুরুত্বপূর্ণ শিক্ষা}

\begin{enumerate}
\item \textbf{তাওবার কেন্দ্র:} তাওবার মূল উপাদান — অন্তরের আফসোস; এই আফসোসই তাকে পাপ ছাড়তে ও সংকল্প করতে শেখায়।
\item \textbf{কপট তাওবা নয়:} শুধু মুখে \lat{“}তাওবা করলাম\lat{”} যথেষ্ট নয় — অন্তরে পাপের জন্য বেদনা চাই; সেই বেদনাই আসল তাওবা।
\item \textbf{উলামাদের ব্যাখ্যা:} হাদিসটি তাওবার \lat{“}মূল\lat{”} বলেছে — বাকি শর্ত (পাপ ছেড়ে দেওয়া, সংকল্প) এই আফসোসের স্বাভাবিক ফল।
\end{enumerate}
\medskip\hrule\medskip

\subsection*{পার্ট ৩ — বড় পাপের পর তাওবা: সাহাবিদের হাদিস-ঘটনা}

\subsubsection*{হাদিস ১৬ — ৯৯ জনের হত্যাকারী — একশোতম হত্যার পরও রহমত (সহীহ মুসলিম ২৭৬৬, সহীহ বুখারি ৩৪৭০)}

\subsubsection*{১. সূত্র কার্ড}

\begin{table}[h]\centering\small
\begin{tabular}{ll}
বিষয় & বিবরণ \\
\hline
\textbf{বর্ণনাকারী} & আবু সাঈদ খুদরি (রাঃ) \\
\textbf{গ্রন্থ} & সহীহ মুসলিম ২৭৬৬; সহীহ বুখারি ৩৪৭০ \\
\textbf{অধ্যায়} & কিতাবুত তাওবাহ — \lat{“}হত্যাকারীর তাওবা কবুল, তার হত্যা যত বেশি হোক\lat{”} \\
\textbf{সনদের মান} & \textbf{সহীহ} (বুখারি-মুসলিম) \\
\hline
\end{tabular}
\end{table}

\subsubsection*{২. পূর্ণ বর্ণনা (বাংলা, \lat{pure})}

\begin{quote}
\textbf{\lat{“}তোমাদের আগে এক ব্যক্তি ছিল, যে নিরানব্বই জনকে হত্যা করেছিল। তারপর সে পৃথিবীর সবচেয়ে জ্ঞানী ব্যক্তির খোঁজ করল। তাকে এক পূজারীর কাছে পাঠানো হলো। সে গিয়ে বলল: আমি নিরানব্বই জনকে হত্যা করেছি — আমার তাওবার সুযোগ আছে কি? পূজারী বলল: না। সে তাকেও হত্যা করল — এতে একশো হয়ে গেল। তারপর আবার জ্ঞানীর খোঁজ করল; তাকে এক আলেমের কাছে পাঠানো হলো। সে বলল: আমি একশো জনকে হত্যা করেছি — আমার তাওবার সুযোগ আছে কি? আলেম বললেন: হ্যাঁ! তাওবার পথে কে বাধা দেবে? তুমি অমুক জায়গায় যাও — সেখানে আল্লাহর ইবাদতকারী লোক আছে; তুমি তাদের সাথে আল্লাহর ইবাদত করো এবং নিজের জায়গায় ফিরে যেও না — কারণ তা মন্দ জায়গা। সে রওনা হলো — অর্ধেক পথ যেতেই তার মৃত্যু এলো। রহমতের ফেরেশতা ও শাস্তির ফেরেশতাদের মধ্যে বিবাদ হলো। রহমতের ফেরেশতারা বলল: সে তাওবাকারী হয়ে আল্লাহর দিকে অন্তর নিয়ে এসেছে। শাস্তির ফেরেশতারা বলল: সে তো কোনো ভালো কাজই করেনি। তখন একজন ফেরেশতা মানুষের রূপে এসে মীমাংসা করলেন: দুই জায়গার দূরত্ব মাপো — যে জায়গার দিকে সে বেশি কাছে, তারই। মাপা হলো — সে ইবাদতের জায়গার বেশি কাছে; রহমতের ফেরেশতারা তাকে নিয়ে গেল।\lat{”}}
\end{quote}

\subsubsection*{৩. শব্দের ব্যাখ্যা}

\begin{table}[h]\centering\small
\begin{tabular}{ll}
শব্দ & অর্থ \\
\hline
\textbf{তিস'আতুন ওয়া তিস'ঈনা নাফসান} & নিরানব্বই জন প্রাণ \\
\textbf{আ'লামি আহলিল আরদ} & পৃথিবীর সবচেয়ে জ্ঞানী \\
\textbf{তাইবুন মুকবিলুন বিক্বলবিহি} & তাওবাকারী — অন্তর নিয়ে আল্লাহর দিকে ফেরা \\
\textbf{ইলাল আরদিল্লাতী আরাদা} & যে জায়গার দিকে সে যেতে চেয়েছিল \\
\hline
\end{tabular}
\end{table}

\subsubsection*{৪. গুরুত্বপূর্ণ শিক্ষা}

\begin{enumerate}
\item \textbf{হত্যার পরেও তাওবা:} ইচ্ছাকৃত হত্যা — ইসলামের সবচেয়ে বড় পাপগুলোর একটি — তারপরও তাওবার দরজা খোলা; অধ্যায়ের নাম: \lat{“}হত্যাকারীর তাওবা কবুল, তার হত্যা যত বেশি হোক\lat{”}।
\item \textbf{আলেমের উত্তর:} \lat{“}তাওবার পথে কে বাধা দেবে?\lat{”} — আপাতদৃষ্টিতে সবচেয়ে বাধা মনে হয় নিজের পাপ; কিন্তু আলেম বললেন — কোনো বাধাই নেই।
\item \textbf{নিয়তের মূল্য:} অন্তরের ফেরাই মূল — সে মৃত্যুর সময় ইবাদতের জায়গার দিকে ছিল; আল্লাহ অন্তর দেখেন।
\item \textbf{পরিবেশের গুরুত্ব:} তাওবার সাথে পরিবেশ বদল — মন্দ জায়গা ছেড়ে ভালো জায়গা; এটাই তাওবার ব্যবহারিক রূপ।
\end{enumerate}
\medskip\hrule\medskip

\subsubsection*{হাদিস ১৭ — বেশ্যা ও তৃষ্ণার্ত কুকুর (সহীহ বুখারি ৩৩২১, সহীহ মুসলিম ২২৪৫)}

\subsubsection*{১. সূত্র কার্ড}

\begin{table}[h]\centering\small
\begin{tabular}{ll}
বিষয় & বিবরণ \\
\hline
\textbf{বর্ণনাকারী} & আবু হুরাইরা (রাঃ) \\
\textbf{গ্রন্থ} & সহীহ বুখারি ৩৩২১; সহীহ মুসলিম ২২৪৫ \\
\textbf{অধ্যায়} & বুখারি: কিতাবু বাদয়িল খালক; মুসলিম: কিতাবুস সালাম — পশুকে পানি দেওয়ার ফজিলত \\
\textbf{সনদের মান} & \textbf{সহীহ} (বুখারি-মুসলিম) \\
\hline
\end{tabular}
\end{table}

\subsubsection*{২. পূর্ণ বর্ণনা (বাংলা, \lat{pure})}

\begin{quote}
\textbf{\lat{“}এক বেশ্যা (পাপী নারী) একটি কুকুরকে কূপের পাশে দেখল — গরমের দিনে সে তৃষ্ণায় জিভ বের করেছিল। সে তার জুতো খুলে তা দিয়ে পানি তুলে কুকুরটিকে পান করাল। এ কারণে আল্লাহ তাকে ক্ষমা করে দিলেন।\lat{”}}
\end{quote}

\subsubsection*{৩. শব্দের ব্যাখ্যা}

\begin{table}[h]\centering\small
\begin{tabular}{ll}
শব্দ & অর্থ \\
\hline
\textbf{বাগিয়্যুন (\arab{بغيّ})} & বেশ্যা — যে নারী যিনায় লিপ্ত \\
\textbf{ইয়ুতীফু (\arab{يطيف})} & চক্কর দিচ্ছিল — কূপের চারপাশে \\
\textbf{মুকাহা (\arab{موقها})} & তার জুতা/মোজা \\
\textbf{ফাগুফিরা লাহা} & তাই তাকে ক্ষমা করা হলো \\
\hline
\end{tabular}
\end{table}

\subsubsection*{৪. গুরুত্বপূর্ণ শিক্ষা}

\begin{enumerate}
\item \textbf{পাপের আকার বনাম রহমত:} জীবনের যিনা-পাপের বিপরীতে একটি পানি পান করানো — রহমত সব পাপ ছাপিয়ে গেল; এটাই হাদিসের মূল বার্তা — ক্ষমার দরজা পাপের হিসাবের চেয়ে বড়।
\item \textbf{ভালো কাজের সুযোগ:} যে অবস্থাতেই থাকো — একটি খাঁটি ভালো কাজ রহমতের দরজা খুলে দিতে পারে।
\item \textbf{পাপীর সম্ভাবনা:} যিনা এমন পাপ যা নবীজির (সা.) যুগে শাস্তির বিষয় ছিল — তবু এক ইখলাসের কারণে ক্ষমা; আল্লাহর রহমত বিচারকের রাগ নয়, পিতার দয়া।
\end{enumerate}
\medskip\hrule\medskip

\subsubsection*{হাদিস ১৮ — মদপানকারী সাহাবি: \lat{“}অভিশাপ দিও না\lat{”} (সহীহ বুখারি ৬৭৮০)}

\subsubsection*{১. সূত্র কার্ড}

\begin{table}[h]\centering\small
\begin{tabular}{ll}
বিষয় & বিবরণ \\
\hline
\textbf{বর্ণনাকারী} & উমার ইবনে খাত্তাব (রাঃ) \\
\textbf{গ্রন্থ} & সহীহ বুখারি ৬৭৮০; সহীহ মুসলিম ২৭৬৮ \\
\textbf{অধ্যায়} & কিতাবুল হুদুদ — মদ্যপায়ীকে অভিশাপ দেওয়া অপছন্দ ও সে কাফির নয় \\
\textbf{সনদের মান} & \textbf{সহীহ} (বুখারি-মুসলিম) \\
\hline
\end{tabular}
\end{table}

\subsubsection*{২. পূর্ণ বর্ণনা (বাংলা, \lat{pure})}

\begin{quote}
নবীজির (সা.) যুগে আবদুল্লাহ নামে এক ব্যক্তি ছিল, যার ডাকনাম ছিল \lat{“}হিমার\lat{”} (গাধা); সে নবীজিকে (সা.) হাসাতেন। মদপানের কারণে নবীজি (সা.) তাকে দণ্ড দিয়েছেন। একদিন তাকে আবার একই অপরাধে আনা হলো এবং দণ্ড দেওয়া হলো। তখন জনতার মধ্যে এক ব্যক্তি বলল: \textbf{\lat{“}হে আল্লাহ! তাকে অভিশাপ দাও — কতবারই না তাকে (এ অপরাধে) আনা হলো!\lat{”}} নবীজি (সা.) বললেন: \textbf{\lat{“}তাকে অভিশাপ দিও না! আল্লাহর কসম! আমি জানি — সে আল্লাহ ও তাঁর রাসূলকে ভালোবাসে।\lat{”}}
\end{quote}

\subsubsection*{৩. শব্দের ব্যাখ্যা}

\begin{table}[h]\centering\small
\begin{tabular}{ll}
শব্দ & অর্থ \\
\hline
\textbf{লাকাবুহু হিমারান} & তার ডাকনাম ছিল হিমার (গাধা) \\
\textbf{জালাদাহু ফিশ শারাব} & মদের জন্য তাকে দণ্ড দিয়েছেন \\
\textbf{লা তাল'আনূহু} & তাকে অভিশাপ দিও না \\
\hline
\end{tabular}
\end{table}

\subsubsection*{৪. গুরুত্বপূর্ণ শিক্ষা}

\begin{enumerate}
\item \textbf{পাপীও মুমিন থাকতে পারে:} বারবার মদপান — বড় পাপ — তবু সে আল্লাহ ও রাসূলকে ভালোবাসে; পাপ ও ঈমান একসাথে থাকতে পারে — তাই পাপীকে \lat{“}হারিয়ে গেছে\lat{”} মনে করা নবীজির (সা.) শিক্ষার বিপরীত।
\item \textbf{পাপীর প্রতি দৃষ্টি:} নবীজি (সা.) পাপকে ঘৃণা করতেন, কিন্তু পাপীকে পুরোপুরি ছেড়ে দেননি — দণ্ড দিয়েছেন, আবার ভালোবাসাও জানতেন।
\item \textbf{অভিশাপ নয়, সংশোধন:} মুসলিমের দায়িত্ব — পাপীর জন্য দোয়া ও সংশোধনের পথ, অভিশাপ নয়; শয়তান চায় পাপীর পাশে কেউ না থাকুক।
\end{enumerate}
\medskip\hrule\medskip

\subsubsection*{হাদিস ১৯ — মায়েয ও গামেদিয়া: পাপ স্বীকার ও ক্ষমার ঘোষণা (সহীহ মুসলিম ১৬৯৫, ১৬৯৬)}

\subsubsection*{১. সূত্র কার্ড}

\begin{table}[h]\centering\small
\begin{tabular}{ll}
বিষয় & বিবরণ \\
\hline
\textbf{বর্ণনাকারী} & বুরাইদা ইবনে হুসাইব (রাঃ); ইমরান ইবনে হুসাইন (রাঃ) \\
\textbf{গ্রন্থ} & সহীহ মুসলিম ১৬৯৫, ১৬৯৬; সহীহ বুখারি ৬৮২৭, ৬৮২৮ \\
\textbf{অধ্যায়} & কিতাবুল হুদুদ — যিনা স্বীকারকারী \\
\textbf{সনদের মান} & \textbf{সহীহ} (বুখারি-মুসলিম) \\
\hline
\end{tabular}
\end{table}

\subsubsection*{২. পূর্ণ বর্ণনা (বাংলা, \lat{pure})}

মায়েয ইবনে মালিক নবীজির (সা.) কাছে এসে বললেন: \textbf{\lat{“}হে আল্লাহর রাসূল! আমাকে পবিত্র করুন।\lat{”}} নবীজি (সা.) বললেন: \textbf{\lat{“}তোমার দুর্ভোগ! ফিরে যাও — আল্লাহর কাছে ক্ষমা চাও এবং তাওবা করো।\lat{”}} সে ফিরে গেল, আবার এলো… চারবার। চতুর্থবার নবীজি (সা.) জিজ্ঞেস করলেন — কিসের থেকে পবিত্র করব? সে বলল — যিনা থেকে। তার পাগলামি ও মদের নেশা যাচাই করা হলো — কিছুই নেই। নবীজি (সা.) তার ব্যাপারে আদেশ দিলেন — তাকে পাথর মারা হলো। তখন লোকেরা দুই দলে ভাগ হলো — কেউ বলল সে ধ্বংস হয়ে গেছে, কেউ বলল — মায়েযের চেয়ে উত্তম তাওবা আর নেই। দুই-তিন দিন পর নবীজি (সা.) সাহাবিদের কাছে এসে বললেন: \textbf{\lat{“}মায়েয ইবনে মালিকের জন্য ক্ষমা চাও।\lat{”}} তারা ক্ষমা চাইলেন। নবীজি (সা.) বললেন: \textbf{\lat{“}সে এমন তাওবা করেছে, যা একটি সম্প্রদায়ের মধ্যে বণ্টন করলে তাদের সকলের জন্য যথেষ্ট হতো।\lat{”}}

এরপর গামেদ (বা জুহাইনা) গোত্রের এক নারী এসে বলল — আমাকে পবিত্র করুন; সে যিনার কারণে গর্ভবতী ছিল। নবীজি (সা.) তাকে সন্তান জন্ম দেওয়া পর্যন্ত অপেক্ষা করতে বললেন; সন্তান হলে তাকে রজম করা হলো; এরপর নবীজি (সা.) তার জানাজার নামাজ পড়ালেন। উমার (রাঃ) বললেন — সে তো যিনা করেছে, আপনি তার জানাজা পড়ালেন? নবীজি (সা.) বললেন: \textbf{\lat{“}সে এমন তাওবা করেছে, যা মদিনার সত্তর জনের মধ্যে বণ্টন করলে তাদের জন্য যথেষ্ট হতো। তুমি কি তার চেয়ে উত্তম তাওবা দেখেছ — যে নিজের প্রাণ আল্লাহর জন্য উৎসর্গ করল?\lat{”}}

\subsubsection*{৩. শব্দের ব্যাখ্যা}

\begin{table}[h]\centering\small
\begin{tabular}{ll}
শব্দ & অর্থ \\
\hline
\textbf{ত্বাহহিরনী (\arab{طهرني})} & আমাকে পবিত্র করুন — শাস্তির মাধ্যমে \\
\textbf{তাবাত তাওবাতান} & এমন তাওবা করেছে \\
\textbf{জাদাত বিনাফসিহা লিল্লাহ} & নিজের প্রাণ আল্লাহর জন্য বিলিয়ে দিয়েছে \\
\hline
\end{tabular}
\end{table}

\subsubsection*{৪. গুরুত্বপূর্ণ শিক্ষা}

\begin{enumerate}
\item \textbf{পাপ স্বীকারের সাহস:} মায়েয ও গামেদিয়া নিজেরা এসে পাপ স্বীকার করেছেন — পাপ ঢাকার সাধারণ নিয়ম (সেটা আলাদা) সত্ত্বেও তাদের স্বীকারোক্তিই তাদের তাওবার সাক্ষ্য।
\item \textbf{তাওবার মান:} নবীজির (সা.) ভাষায় — সম্প্রদায়/সত্তর জনের জন্য যথেষ্ট এমন তাওবা; অর্থাৎ অন্তরের আন্তরিক ফেরা সব পাপের হিসাব মুছে দেয়।
\item \textbf{শাস্তি মানে ধ্বংস নয়:} দুনিয়ার শাস্তি (হদ) আদায় — আখিরাতের শাস্তি থেকে মুক্তি; মায়েযের পেছনে নবীজির (সা.) ক্ষমা প্রার্থনা এর প্রমাণ।
\item \textbf{লোকের বিচার নয়:} সাহাবিরাও দ্বিধাবিভক্ত ছিলেন — কিন্তু আল্লাহ ও তাঁর রাসূলের সিদ্ধান্তই চূড়ান্ত; পাপীর ভবিষ্যৎ আল্লাহর হাতে।
\end{enumerate}
\medskip\hrule\medskip

\subsubsection*{হাদিস ১৯খ — যে যুবক যিনার অনুমতি চেয়েছিল — নবীজির (সা.) দোয়ায় হৃদয় পরিশুদ্ধ (মুসনাদ আহমাদ ২২২১৩)}

\subsubsection*{১. সূত্র কার্ড}

\begin{table}[h]\centering\small
\begin{tabular}{ll}
বিষয় & বিবরণ \\
\hline
\textbf{বর্ণনাকারী} & আবু উমামা বাহিলি (রাঃ) \\
\textbf{গ্রন্থ} & মুসনাদ আহমাদ ২২২১৩; তাবারানির আল-মু'জামুল কবীর; আল-মুগনী আন হামলিল আসফার ২২৬৬ \\
\textbf{অধ্যায়} & মুসনাদ আহমাদ — আবু উমামার বর্ণনা \\
\textbf{সনদের মান} & \textbf{সহীহ/হাসান সনদ} — ইরাকি: \lat{“}সহীহের রাবিদের সনদ\lat{”}; আলবানি: সহীহ \\
\hline
\end{tabular}
\end{table}

\subsubsection*{২. পটভূমি}

এক যুবক নবীজির (সা.) কাছে এসে সরাসরি বলল — আমাকে যিনার অনুমতি দিন। আশেপাশের লোকেরা তাকে ধমক দিতে লাগল। নবীজি (সা.) তাকে কাছে ডেকে বসালেন এবং প্রশ্ন দিয়ে তার অন্তরকে জাগিয়ে দিলেন।

\subsubsection*{৩. পূর্ণ বর্ণনা (বাংলা, \lat{pure})}

\begin{quote}
এক যুবক নবীজির (সা.) কাছে এসে বলল: \textbf{\lat{“}হে আল্লাহর রাসূল! আমাকে যিনার অনুমতি দিন।\lat{”}} লোকেরা তাকে ধমক দিয়ে বলল: থামো, থামো! নবীজি (সা.) বললেন: \textbf{\lat{“}তাকে কাছে আসতে দাও।\lat{”}} সে কাছে এসে বসল। নবীজি (সা.) জিজ্ঞেস করলেন: \textbf{\lat{“}তুমি কি এটা তোমার মায়ের জন্য পছন্দ করবে?\lat{”}} সে বলল: না — আল্লাহর কসম, আপনার জন্য আমি উৎসর্গিত! নবীজি (সা.) বললেন: \textbf{\lat{“}লোকেরাও তাদের মায়ের জন্য তা পছন্দ করে না।\lat{”}} একইভাবে কন্যা, বোন, ফুফু, খালার কথা জিজ্ঞেস করলেন — প্রতিবারই সে বলল না; নবীজি (সা.) বললেন — \textbf{\lat{“}লোকেরাও তাদের জন্য তা পছন্দ করে না।\lat{”}} এরপর নবীজি (সা.) তার উপর হাত রেখে দোয়া করলেন: \textbf{\lat{“}হে আল্লাহ! তার গুনাহ ক্ষমা করো, তার অন্তর পবিত্র করো এবং তার লজ্জাস্থানকে (পাপ থেকে) হিফাজত করো।\lat{”}} বর্ণনাকারী বলেন — এরপর সে যুবক আর কোনো নিষিদ্ধ জিনিসের দিকে ভ্রুক্ষেপ করেনি।
\end{quote}

\subsubsection*{৪. শব্দের ব্যাখ্যা}

\begin{table}[h]\centering\small
\begin{tabular}{ll}
শব্দ & অর্থ \\
\hline
\textbf{আযান লি বিজ যিনা} & আমাকে যিনার অনুমতি দাও \\
\textbf{তাহহির ক্বলবাহু (\arab{طهر} \arab{قلبه})} & তার অন্তর পবিত্র করো \\
\textbf{হাসসিন ফারজাহু (\arab{حصّن} \arab{فرجه})} & তার লজ্জাস্থানকে হিফাজত করো \\
\hline
\end{tabular}
\end{table}

\subsubsection*{৫. গুরুত্বপূর্ণ শিক্ষা}

\begin{enumerate}
\item \textbf{পাপীর জন্য নবীজির (সা.) পদ্ধতি:} ধমক নয় — কাছে ডেকে প্রশ্ন দিয়ে চেতনা জাগানো; পাপীকে ছেড়ে না দেওয়া।
\item \textbf{দোয়ার শক্তি:} নবীজির (সা.) দোয়া — \lat{“}অন্তর পবিত্র করো, লজ্জাস্থান হিফাজত করো\lat{”} — আল্লাহর সাহায্য পাপীকে বদলে দিল; বর্ণনাকারীর ভাষায় সে আর পাপের দিকে তাকায়নি।
\item \textbf{আল্লাহর সাহায্যের নমুনা:} যে সত্যিই ফিরতে চায়, আল্লাহ তাকে পথ দেখান (আনকাবুত ২৯:৬৯) — এই ঘটনা তার জীবন্ত উদাহরণ।
\end{enumerate}
\medskip\hrule\medskip

\subsection*{পার্ট ৪ — তাওবার পরের আমল (\lat{Life} \lat{After} \lat{Tawbah})}

\subsubsection*{হাদিস ২০ — \lat{“}মন্দ কাজের পর ভালো কাজ করো\lat{”} (তিরমিযী ১৯৮৭)}

\subsubsection*{১. সূত্র কার্ড}

\begin{table}[h]\centering\small
\begin{tabular}{ll}
বিষয় & বিবরণ \\
\hline
\textbf{বর্ণনাকারী} & আবু যার গিফারী (রাঃ) \\
\textbf{গ্রন্থ} & জামি' আত-তিরমিযী ১৯৮৭; মুসনাদ আহমাদ \\
\textbf{অধ্যায়} & কিতাবুল বিরর — নেকি ও গুনাহ \\
\textbf{সনদের মান} & \textbf{হাসান-সহীহ} (তিরমিযী) \\
\hline
\end{tabular}
\end{table}

\subsubsection*{২. পূর্ণ বর্ণনা (বাংলা, \lat{pure})}

\begin{quote}
\textbf{\lat{“}যেখানে (যে অবস্থায়) তুমি থাকো, আল্লাহকে ভয় করো; মন্দ কাজের পর ভালো কাজ করো — তা তা (মন্দ কাজ) মুছে দেবে; আর মানুষের সাথে সুন্দর আচরণ করো।\lat{”}}
\end{quote}

\subsubsection*{৩. গুরুত্বপূর্ণ শিক্ষা}

\begin{enumerate}
\item \textbf{তাওবার বাস্তব রূপ:} পাপের পর ভালো কাজ — তাওবার প্রমাণ ও পাপ মোছার মাধ্যম; নিছক আফসোস নয়, আমল।
\item \textbf{ভয় ও আশার সংমিশ্রণ:} \lat{“}আল্লাহকে ভয় করো\lat{”} — পাপ বন্ধের ভয়; \lat{“}ভালো কাজ করো\lat{”} — ক্ষমার আশা; দুই-ই এক হাদিসে।
\end{enumerate}
\medskip\hrule\medskip

\subsubsection*{হাদিস ২১ — সাদাকা পাপ নিভিয়ে দেয় (তিরমিযী ২৬১৬)}

\subsubsection*{১. সূত্র কার্ড}

\begin{table}[h]\centering\small
\begin{tabular}{ll}
বিষয় & বিবরণ \\
\hline
\textbf{বর্ণনাকারী} & মুআয ইবনে জাবাল (রাঃ) \\
\textbf{গ্রন্থ} & জামি' আত-তিরমিযী ২৬১৬; ইবনে মাজাহ ৩৯৭৩ \\
\textbf{অধ্যায়} & কিতাবুল জুমুআ — সাদাকার ফজিলত \\
\textbf{সনদের মান} & \textbf{হাসান} (দারুসসালাম: হাসান) \\
\hline
\end{tabular}
\end{table}

\subsubsection*{২. পূর্ণ বর্ণনা (বাংলা, \lat{pure})}

\begin{quote}
\textbf{\lat{“}সাদাকা পাপকে নিভিয়ে দেয়, যেমন পানি আগুনকে নিভিয়ে দেয়।\lat{”}}
\end{quote}

\subsubsection*{৩. গুরুত্বপূর্ণ শিক্ষা}

\begin{enumerate}
\item \textbf{পাপের প্রতিষেধক:} তাওবার সাথে দান-সাদাকা — পাপের আগুন নেভানোর অন্যতম উপায়; তাওবাকারী পাপীর জন্য ব্যবহারিক পথ।
\item \textbf{উপমা:} পানি যেমন আগুন নেভায় — সাদাকা তেমনি পাপ মুছে দেয়; তাওবার পরে দান তাওবার সীল।
\end{enumerate}
\medskip\hrule\medskip

\subsubsection*{হাদিস ২২ — ইউনুসের (আঃ) দোয়া — সাড়া পাওয়ার দোয়া (তিরমিযী ৩৫০৫)}

\subsubsection*{১. সূত্র কার্ড}

\begin{table}[h]\centering\small
\begin{tabular}{ll}
বিষয় & বিবরণ \\
\hline
\textbf{বর্ণনাকারী} & সা'দ ইবনে আবি ওয়াক্কাস (রাঃ) \\
\textbf{গ্রন্থ} & জামি' আত-তিরমিযী ৩৫০৫ \\
\textbf{অধ্যায়} & কিতাবুদ দাওয়াত — যুন-নূনের দোয়া \\
\textbf{সনদের মান} & \textbf{হাসান-সহীহ} (তিরমিযী; আলবানি: সহীহ) \\
\hline
\end{tabular}
\end{table}

\subsubsection*{২. পূর্ণ বর্ণনা (বাংলা, \lat{pure})}

\begin{quote}
\textbf{\lat{“}যুন-নূন (ইউনুস আঃ)-এর দোয়া — যখন সে মাছের পেটে দোয়া করেছিল: লা ইলাহা ইল্লা আনতা সুবহানাকা ইন্নি কুনতু মিনাজ জালিমিন (\lat{“}তুমি ছাড়া কোনো সত্য ইলাহ নেই — তুমি পবিত্র; আমি জালিমদের অন্তর্ভুক্ত\lat{”}) — কোনো মুসলিম ব্যক্তি এটি দিয়ে কোনো বিষয়ে দোয়া করলে আল্লাহ তার দোয়ার জবাব দেন না — এমনটি হয় না।\lat{”}}
\end{quote}

\subsubsection*{৩. শব্দের ব্যাখ্যা}

\begin{table}[h]\centering\small
\begin{tabular}{ll}
শব্দ & অর্থ \\
\hline
\textbf{যুন-নূন (\arab{ذو} \arab{النون})} & মাছওয়ালা — ইউনুস (আঃ) \\
\textbf{ইন্নি কুনতু মিনাজ জালিমিন} & আমি জালিমদের অন্তর্ভুক্ত — দোষ স্বীকার \\
\hline
\end{tabular}
\end{table}

\subsubsection*{৪. গুরুত্বপূর্ণ শিক্ষা}

\begin{enumerate}
\item \textbf{পাপীর দোয়া:} দোষ স্বীকারের দোয়া — আল্লাহর কাছে পৌঁছানোর সবচেয়ে কার্যকর পথ; নবীও এই দোয়ায় ফিরেছেন।
\item \textbf{সাড়ার প্রতিশ্রুতি:} যত বড় বিপদ/পাপই হোক — এই দোয়ার মাধ্যমে আল্লাহর সাড়ার নিশ্চয়তা; তাওবাকারীর জন্য প্রস্তুত দোয়া।
\end{enumerate}
\medskip\hrule\medskip

\subsubsection*{হাদিস ২৩ — ঋণদাতার ক্ষমা — আল্লাহর ক্ষমার কারণ (সহীহ বুখারি ২০৭৮, ৩৪৮০; সহীহ মুসলিম ১৫৬২)}

\subsubsection*{১. সূত্র কার্ড}

\begin{table}[h]\centering\small
\begin{tabular}{ll}
বিষয় & বিবরণ \\
\hline
\textbf{বর্ণনাকারী} & আবু হুরাইরা (রাঃ) \\
\textbf{গ্রন্থ} & সহীহ বুখারি ২০৭৮, ৩৪৮০; সহীহ মুসলিম ১৫৬২ \\
\textbf{অধ্যায়} & কিতাবুল বুয়ূ' — ঋণ সংক্রান্ত; মুসলিম: কিতাবুল মুসাকাত \\
\textbf{সনদের মান} & \textbf{সহীহ} (বুখারি-মুসলিম) \\
\hline
\end{tabular}
\end{table}

\subsubsection*{২. পূর্ণ বর্ণনা (বাংলা, \lat{pure})}

\begin{quote}
\textbf{\lat{“}তোমাদের আগের এক সম্প্রদায়ের এক ব্যক্তি ছিল, যে মানুষকে ঋণ দিত। সে তার কর্মচারীকে বলত: কারো কাছে গেলে, সে যদি অসচ্ছল হয়, তবে তাকে ছেড়ে দাও (মাফ করে দাও) — আশা করি আল্লাহও আমাদের ক্ষমা করবেন। আল্লাহ তাকে ক্ষমা করে দিলেন।\lat{”}}
\end{quote}

\subsubsection*{৩. গুরুত্বপূর্ণ শিক্ষা}

\begin{enumerate}
\item \textbf{ক্ষমা পেতে ক্ষমা:} মানুষকে ক্ষমা করলে আল্লাহ ক্ষমা করেন — তাওবাকারী পাপীর জন্য অন্যতম মাধ্যম; নিজের পাপের ক্ষমা চাইলে অন্যদের ত্রুটি ক্ষমা করা তার সঙ্গী হওয়া উচিত।
\item \textbf{নিয়তের শর্ত:} তার নিয়ত ছিল — \lat{“}আশা করি আল্লাহ আমাদের ক্ষমা করবেন\lat{”} — মানুষের প্রতি দয়া ও আল্লাহর প্রতি আশা একসাথে।
\end{enumerate}
\medskip\hrule\medskip

\subsubsection*{হাদিস ২৪ — গুহায় তিনজন: পাপ এড়ানোর দোয়া ও আল্লাহর সাহায্য (সহীহ মুসলিম ২৭৪৩, সহীহ বুখারি ২২১৫)}

\subsubsection*{১. সূত্র কার্ড}

\begin{table}[h]\centering\small
\begin{tabular}{ll}
বিষয় & বিবরণ \\
\hline
\textbf{বর্ণনাকারী} & আবদুল্লাহ ইবনে উমার (রাঃ) \\
\textbf{গ্রন্থ} & সহীহ মুসলিম ২৭৪৩; সহীহ বুখারি ২২১৫ \\
\textbf{অধ্যায়} & মুসলিম: কিতাবুর রিকাক — গুহার তিনজনের ঘটনা \\
\textbf{সনদের মান} & \textbf{সহীহ} (বুখারি-মুসলিম) \\
\hline
\end{tabular}
\end{table}

\subsubsection*{২. পূর্ণ বর্ণনা (বাংলা, \lat{pure})}

\begin{quote}
তিন ব্যক্তি সফরে গুহায় আশ্রয় নিলেন — পাহাড়ের পাথর মুখ বন্ধ করে দিল। তারা বলল — আল্লাহর সন্তুষ্টির জন্য করা ভালো আমলের উসিলায় দোয়া করি। দ্বিতীয়জন বলল: \textbf{\lat{“}হে আল্লাহ! আমার এক ফুফাতো বোন ছিল — তাকে আমি পুরুষের মতো ভালোবাসতাম; আমি তার কাছে চাইলাম, সে একশো দিনার ছাড়া রাজি হলো না। কষ্ট করে একশো দিনার জমা করে তার কাছে দিলাম। যখন তার উপর বসতে যাচ্ছিলাম, সে বলল: আল্লাহকে ভয় করো — বৈধ পথ ছাড়া সীল ভাঙো না। আমি উঠে গেলাম, তাকে একশো দিনার রেখে দিলাম। হে আল্লাহ! আমি যদি এটা শুধু তোমার সন্তুষ্টির জন্য করে থাকি, তবে আমাদের এই বিপদ দূর করো।\lat{”} পাথর কিছুটা সরে গেল…}
\end{quote}

\subsubsection*{৩. শব্দের ব্যাখ্যা}

\begin{table}[h]\centering\small
\begin{tabular}{ll}
শব্দ & অর্থ \\
\hline
\textbf{ইবনাতু আম্মি (\arab{ابنة} \arab{عمّي})} & ফুফাতো বোন \\
\textbf{সিকরুন (\arab{سكرا})} & সীল/বন্ধন — এখানে কুমারীত্বের রূপক \\
\textbf{ইত্তাকিল্লাহ} & আল্লাহকে ভয় করো — তার উপদেশ \\
\hline
\end{tabular}
\end{table}

\subsubsection*{৪. গুরুত্বপূর্ণ শিক্ষা}

\begin{enumerate}
\item \textbf{পাপের শেষ মুহূর্তে ফেরা:} সবকিছু প্রস্তুত — টাকাও দেওয়া হয়েছে — তবু \lat{“}আল্লাহকে ভয় করো\lat{”} শুনে পাপ ছেড়ে দেওয়া; এই ফেরাই আল্লাহর সাহায্য টানে।
\item \textbf{আল্লাহর সাহায্য:} পাপ এড়ানোর ইখলাসের কারণে আল্লাহ বিপদ থেকে উদ্ধার করলেন — তাওবাকারীর জন্য আল্লাহর সাহায্যের নমুনা।
\end{enumerate}
\medskip\hrule\medskip

\subsection*{পার্ট ৫ — যাচাই টেবিল: প্রচলিত-কিন্তু-দুর্বল/ভিত্তিহীন বিষয় (\lat{Verification} \lat{Table})}

\begin{quote}
\textbf{সতর্কতা:} নিচের বিষয়গুলো প্রচলিত, কিন্তু \textbf{\lat{authentic} নয়}। এগুলো দলিল হিসেবে ব্যবহার করা যাবে না। \lat{AGENTS}.\lat{md}-এর নিয়ম অনুযায়ী এগুলো মূল সংকলনে আনা হয়নি; শুধু যেন কেউ এগুলোকে সহীহ ভেবে না নেয়, সেজন্য এখানে স্পষ্টভাবে চিহ্নিত করা হলো।
\end{quote}

\begin{table}[h]\centering\small
\begin{tabular}{llll}
\# & প্রচলিত বিষয় & আসল অবস্থা & ব্যাখ্যা ও উৎস \\
\hline
১ & ফুদাইল ইবনে ইয়াদের (রহ.) ডাকাত-জীবন থেকে তাওবার গল্প (রাতে কাফেলার লোকেরা তার নামে ভয় পেলে তাওবা করেন) & \textbf{দুর্বল সনদ (যয়ীফ)} & সিয়ারু আলাম আন-নুবালা (৮/৪৩৮), তারিখু দিমাশক (৪৮/৩৮৪)-এর সনদ দুর্বল — অজ্ঞাত/দুর্বল রাবি, পরোক্ষ সূত্র। ফুদাইল (রহ.) নিজে একজন বড় আলেম ছিলেন — এটা ঐতিহাসিক সত্য; কিন্তু ডাকাতির এই বিশেষ গল্পটির সনদ বিশ্বাসযোগ্য নয়। \\
২ & \lat{“}তুমি গুনাহের দিকে তাকিও না — তাকাও সেই আল্লাহর দিকে, যাঁর কাছে তুমি গুনাহ করেছ\lat{”} & \textbf{হাদিস নয় — উলামাদের উক্তি} & এটি অনেক সময় নবীজির (সা.) নামে প্রচারিত হয় — ভুল। এটি আলেমদের (ইবনে কাইয়্যিম, ইবনে আতা ইল্লাহ প্রমুখের) উক্তির সারমর্ম; হাদিস নয়। \\
৩ & \lat{“}শুধু মুখে 'আমি তাওবা করলাম' বললেই সব মাফ — আমল লাগবে না\lat{”} & \textbf{ভ্রান্ত ধারণা (\lat{misconception})} & তাওবার শর্ত — পাপ ছেড়ে দেওয়া, আফসোস, সংকল্প (তিরমিযী ১৯৮৭-এর মতো আমলের প্রমাণ); মুখে বলা যথেষ্ট নয়। ৬৬:৮-তে তাওবাতুন নসূহের আদেশ। \\
৪ & \lat{“}অতীতের পাপের কারণে আল্লাহ আর কোনো দোয়া কবুল করেন না\lat{”} & \textbf{ভিত্তিহীন} & কুরআন-হাদিসে এমন কোনো দলিল নেই; বরং ২:১৮৬-তে ডাকার জবাবের প্রতিশ্রুতি, ৪০:৬০-তে দোয়ার প্রতিশ্রুতি — পাপীকে ডাকা থেকে বিরত রাখার দলিল কোথাও নেই। \\
৫ & \lat{“}এক কান্না/একবার ইস্তিগফারেই সব পাপের হিসাব শেষ — আর পাপ করলেই হবে না\lat{”} & \textbf{ভুল বোঝাবুঝি} & ক্ষমার দরজা খোলা (মুসলিম ২৭৫৯) মানে পাপে নির্ভীক হওয়া নয়; বরং খাঁটি তাওবা (৬৬:৮) ও পাপ ছেড়ে দেওয়া (৩:১৩৫) শর্ত — আশা ও আমল দুই-ই। \\
৬ & ইবরাহীম ইবনে আদহাম, মালিক ইবনে দিনার প্রমুখের রাজকীয়/ডাকাত জীবন থেকে তাওবার বিস্তারিত গল্প & \textbf{ঐতিহাসিক বর্ণনা — হাদিস-গ্রেড নয়} & এসব সীরাত/আখবার (ইতিহাস) ও তাযকিয়া সাহিত্যের (হিলয়াতুল আউলিয়া প্রভৃতি) বর্ণনা; সনদ সহীহ হাদিসের মানে নয়। শিক্ষা গ্রহণ করা যায়, কিন্তু \lat{“}সহীহ হাদিসের ঘটনা\lat{”} হিসেবে বলা যাবে না। \\
৭ & \lat{“}আল্লাহ বলেন: হে আদম সন্তান, আমি অসুস্থ ছিলাম, তুমি আমাকে দেখতে এলে না…\lat{”} — এটি দিয়ে \lat{“}আল্লাহর পক্ষ থেকে পাপীর ক্ষমার গল্প\lat{”} বানানো & \textbf{হাদিসটির ব্যবহার ভুল} & হাদিসটি সহীহ (মুসলিম ২৫৬৯) — তবে এটি পাপ-ক্ষমার নয়; এটি অন্য মুসলিমের হকের কথা (অসুস্থ, ক্ষুধার্ত, বন্দিকে সাহায্য — আল্লাহর নামে)। অর্থ বিকৃত করে অন্য প্রসঙ্গে ব্যবহার করা যাবে না। \\
৮ & ওয়াহশি ইবনে হারব (হামযার (রা.) হত্যাকারী)-এর তাওবার গল্প — \lat{“}হামযাকে হত্যার পর নবীজির (সা.) কাছে ফিরে আসা, ২৫:৭০ নাজিল, শেষে মুসাইলামাকে হত্যা\lat{”} & \textbf{সীরাত-গ্রেড (ঐতিহাসিক) — সহীহ হাদিস নয়} & ঘটনাটি ইবনে ইসহাকের সীরাত ও তাফসীর-ইতিহাসের (তাবারী, ইবনে কাসীর) বর্ণনা; সহীহ হাদিস গ্রন্থে মারফু' সনদসহ নেই। ২৫:৭০-এর শানে নুজুল হিসেবে ওয়াহশির ঘটনাটিও দুর্বল সনদে বর্ণিত — ইবনে কাসীর নিজেই এ নিয়ে সতর্ক করেছেন। শিক্ষাগত মূল্য আছে, কিন্তু \lat{“}সহীহ হাদিসের ঘটনা\lat{”} বলা যাবে না। \\
৯ & আবু লুবাবার (রা.) স্তম্ভে নিজেকে বেঁধে রাখার ঘটনা (বনু কুরাইজার ঘটনায়) — ৮:২৭ নাজিল পর্যন্ত & \textbf{তাফসীর/সীরাত-গ্রেড — সহীহ হাদিস সনদসহ নেই} & গল্পটি ইবনে ইসহাকের সীরাত ও তাবারী/ইবনে কাসীরের তাফসীরে আছে; সহীহ বুখারি-মুসলিমে এই ঘটনাটি সনদসহ নেই। নবীজির (সা.) জীবনের ঐতিহাসিক ঘটনা হিসেবে আলোচনা করা যায় — কিন্তু সহীহ হাদিস হিসেবে নয়। \\
১০ & থুমামাহ ইবনে উথালের ঘটনা — বন্দি হওয়া, নবীজি (সা.)-এর মুক্তি ও ইসলাম গ্রহণ & \textbf{সীরাত-গ্রেড (ঐতিহাসিক)} & ইবনে ইসহাকের সীরাতে আছে; সহীহ হাদিস গ্রন্থে মারফু' সনদসহ নেই। ঐতিহাসিক ঘটনা — কিন্তু এখানে ব্যবহৃত দলিল-মানের (সহীহ হাদিস) তালিকায় রাখা যায় না। \\
১১ & বিশর আল-হাফির (রহ.) তাওবার গল্প — রাস্তায় \lat{“}বিসমিল্লাহ\lat{”} লেখা কাগজ তুলে সুগন্ধি দেওয়া & \textbf{দুর্বল/পরস্পরবিরোধী সনদ} & ঘটনাটি তাযকিয়া সাহিত্যে (সিয়ার, হিলয়াহ) প্রচলিত; কিন্তু সনদ দুর্বল বা পরস্পরবিরোধী — কঠোর মানে সহীহ বলা যায় না। বিশর (রহ.) নিজে একজন বড় আলেম-আবিদ ছিলেন (এটা ঐতিহাসিক সত্য); কিন্তু তার তাওবার এই বিশেষ গল্পটির সনদ বিশ্বাসযোগ্য নয়। \\
\hline
\end{tabular}
\end{table}

\begin{quote}
\textbf{নিয়ম:} উপরের কোনো বিষয়কে এই সংকলনের \lat{“}ঘটনা\lat{”} অংশে \lat{authentic} ঘটনা হিসেবে আনা হয়নি। সাহাবিদের ঘটনা ফাইলে কেবল সহীহ (বুখারি-মুসলিম বা হাসান) সূত্রের ঘটনা আছে।
\end{quote}

\medskip\hrule\medskip

\subsection*{রেফারেন্স তালিকা}

\begin{itemize}
\item সহীহ বুখারি: ৩১৯৪, ৩৩২১, ৩৪৭০, ৩৪৭৮, ৫৬৭৩, ৫৯৯৯, ৬০০০, ৬৩০৮, ৬৭৮০, ৬৮২৭-৬৮২৮, ৭৪০৫, ৭৪৫৩, ৭৫০৭, ২০৭৮, ৩৪৮০, ২২১৫
\item সহীহ মুসলিম: ২২৪৫, ২৫৬৯, ২৫৭৭, ২৬৭৫, ২৭৪৩, ২৭৪৪, ২৭৪৯, ২৭৫১, ২৭৫২, ২৭৫৪, ২৭৫৬, ২৭৫৮, ২৭৫৯, ২৭৬৬, ২৭৬৮, ২৮১৬, ১৬৯৫, ১৬৯৬, ১৫৬২
\item জামি' আত-তিরমিযী: ১৯৮৭ (হাসান-সহীহ), ২৪৯৯ (দুর্বল সনদ — মতভেদ), ২৬১৬ (হাসান), ৩৫৩৭ (হাসান), ৩৫৪০ (হাসান), ৩৫০৫ (হাসান-সহীহ)
\item সুনানে ইবনে মাজাহ: ৪২৫১ (হাসান), ৪২৫২ (হাসান), ৪২৫৭, ৩৯৭৩
\item মুসনাদ আহমাদ: ২২২১৩ (যিনা-অনুমতি চাওয়া যুবকের হাদিস — ইরাকি ও আলবানি: সনদ সহীহ/হাসান)
\item ব্যাখ্যা: ফাতহুল বারি (ইবনে হাজার), শারহু নববী (সহীহ মুসলিম), তাফসীর ইবনে কাসীর
\end{itemize}

\end{document}
